% ============================================================================
% Chapter 2, Section 4: Formalization
% ============================================================================

\subsection{Formalization}\label{sec:formalization}

We establish the formal framework used throughout the remainder of this thesis. The definitions below are adapted from the MetaHarness formalism~\citep{karten2025pokeagent} and extended to accommodate the harness spectrum that our experiments explore.

\subsubsection{Embodied Agent Environments}\label{sec:formal_env}

We consider an embodied agent interacting with a visual environment through a minimal interface. At each discrete timestep $t$, the agent receives a frame observation $o_t$ (a rendered image of the current game state) and a structured text state $x_t$ (derived from read-only memory and map data), then selects an action $a_t$ from a fixed set of button inputs $\mathcal{A}$. The environment is partially observable: the agent cannot access the full internal game state---such as off-screen NPC positions, hidden item locations, or opponent battle mechanics---beyond what is visually rendered or explicitly provided through the text state.

Formally, we model the environment as a partially observable process:
\begin{equation}\label{eq:env}
    (o_{t+1}, x_{t+1}) = \mathcal{E}(s_t, a_t), \quad s_{t+1} = f(s_t, a_t),
\end{equation}
where $s_t \in \mathcal{S}$ is the hidden game state, $f$ is the deterministic (up to random encounters) state transition function implemented by the game engine, and $\mathcal{E}$ is the observation function that renders the visible portion of the state. This formulation is general: any environment that exposes visual output and accepts discrete inputs fits this interface, making our framework applicable beyond Pok\'{e}mon to other visual, embodied domains.

The agent maintains a trajectory $\tau_t = \{(o_1, x_1, a_1), \ldots, (o_t, x_t, a_t)\}$ recording its full interaction history. In practice, the trajectory also includes the agent's reasoning traces (chain-of-thought outputs) and tool-call records at each step, forming a rich log that serves as the substrate for both context management and evolutionary optimization.

\subsubsection{Agentic Harnesses}\label{sec:formal_harness}

An \emph{agentic harness} $\mathcal{H}$ is the scaffolding layer between a foundation model $M$ and the environment. Following the decomposition introduced in the \pokeagent{} Challenge~\citep{karten2025pokeagent}, a harness mediates agent behavior through five components:
\begin{itemize}[leftmargin=*]
    \item \textbf{System prompt} $p$: instructions and strategic guidance provided to the model at each reasoning step.
    \item \textbf{Subagents} $\mathcal{S} = \{s_1, \ldots, s_k\}$: specialized modules that can be invoked by the orchestrator for specific tasks (e.g., battle strategy, navigation planning, self-reflection).
    \item \textbf{Memory} $\mathcal{M}$: a persistent knowledge store that accumulates facts, strategies, and observations across the agent's trajectory.
    \item \textbf{Skills} $\mathcal{K}$: reusable behavioral routines---either textual (strategic guidance) or executable (Python code in a sandbox)---that encode solutions to previously encountered situations.
    \item \textbf{Tools} $\mathcal{T}$: programmatic capabilities available to the model, such as pathfinding, information lookup, or code execution.
\end{itemize}

We define a spectrum of harness configurations that characterize different levels of scaffolding:

\begin{description}[leftmargin=*]
    \item[$\mathcal{H}_{\min}$ (Minimalist):] Provides only the environment interface---frame observations and button inputs---with a generic system prompt and no subagents, memory, skills, or domain-specific tools. This represents the lower bound on scaffolding: the model must rely entirely on its pretrained knowledge and in-context reasoning.

    \item[$\mathcal{H}_{\text{expert}}$ (Human-Designed):] All components are populated through manual engineering. In our setting, this is the full PokeAgent harness described in Chapter~4, with hand-crafted subagents for reflection, battle strategy, and navigation planning; a curated knowledge base; walkthrough-derived objectives; and Porymap-based~\citep{huderlem2025porymap} spatial tools.

    \item[$\mathcal{H}_{\text{meta}}$ (Meta-Tool):] An intermediate approach that gives the model meta-tools for constructing its own agents and tools (e.g., \texttt{define\_agent}, \texttt{run\_code}, custom tool creation) rather than pre-building specific components. This was the approach taken in later Gemini Plays Pok\'{e}mon runs~\citep{zhang2025gpp}, where the model spontaneously constructed pathfinders, battle strategists, and reusable scripts.

    \item[$\mathcal{H}_{\text{auto}}$ (Auto-Evolved):] Starts from $\mathcal{H}_{\min}$ augmented with meta-tools and adds an evolutionary optimization loop that automatically triggers harness refinement based on trajectory analysis. This is the AutoEvolve output, which evolves toward a harness that may approach---or differ qualitatively from---$\mathcal{H}_{\text{expert}}$ during a single episode.
\end{description}

The central question of Chapter~5 is whether $\mathcal{H}_{\text{auto}}$, evolved from $\mathcal{H}_{\min}$ through AutoEvolve's reset-free optimization, can approach the performance of $\mathcal{H}_{\text{expert}}$---and if so, what the resulting harness looks like compared to the human-designed version.

\subsubsection{In-Context Prompt Evolution}\label{sec:formal_evolution}

Prompt evolution methods treat the system prompt as a parameter optimized using trajectory outcomes as signal. Let $P_i$ denote the system prompt at evolution step $i$, and let $\Delta\tau_i = \{\tau_{t-w}, \ldots, \tau_t\}$ denote a trajectory segment of window size $w$ centered at the current timestep.

In the \emph{episodic} setting (e.g., GEPA~\citep{agrawal2025gepa}), optimization proceeds across complete episodes:
\begin{equation}\label{eq:episodic}
    P_{i+1} = g_{\text{opt}}\!\left(P_i,\; f_{\text{judge}}(\tau^{(i)}_{\text{full}})\right),
\end{equation}
where $\tau^{(i)}_{\text{full}}$ is the complete trajectory from episode $i$, $f_{\text{judge}}$ produces a structured evaluation, and $g_{\text{opt}}$ synthesizes an updated prompt. Each iteration requires resetting the environment to the initial state.

In the \emph{reset-free} setting that AutoEvolve adopts, evolution occurs mid-episode using trajectory segments:
\begin{equation}\label{eq:resetfree}
    \mathcal{H}_{i+1} = g_{\text{evolve}}\!\left(\mathcal{H}_i,\; f_{\text{judge}}(\Delta\tau_i)\right),
\end{equation}
where the evolution operator $g_{\text{evolve}}$ modifies the full harness $\mathcal{H}_i = (P_i, \mathcal{S}_i, \mathcal{K}_i, \mathcal{M}_i)$---not just the prompt---based on the judge's analysis of the recent trajectory segment. This formulation enables continuous adaptation without interrupting the agent's progress and generalizes episodic prompt evolution (Equation~\ref{eq:episodic}) to the case of partial trajectories and multi-component harness optimization. The detailed mechanism of $g_{\text{evolve}}$, including its four-pass structure over prompts, subagents, skills, and memory, is presented in Chapter~5.
